\begin{pastexam}[label={px:ntu-108-dsaa-06}]{臺大資工所 108 學年度 第 6 題}
（問答）

(13\%) When implementing a hash table, one of the commonly used collision resolution methods is
open addressing. Let $k$ be the key, and $m$ be the size of the hash table. Assume each table entry
can hold one key.
\begin{enumerate}[label=(\alph*)]
\item (3\%) The simplest open addressing method is called \textbf{linear probing}, where if a slot
is already occupied, the following slots are ``probed'' sequentially. Assume an auxiliary hash
function $h'(k)=k \bmod m$. Please give the hash function $h(k,i)$ of linear probing with the given
auxiliary function, where $i$ is the number of times the hash table has been probed with this key
(and thus starting from 0).
\item (5\%) Linear probing is easy to implement but suffers from the problem of primary
clustering. A more complex open addressing method which can mitigate this problem is
\textbf{double hashing}, using a hash function in the form of
$h(k,i)=\big(h_1(k)+ih_2(k)\big) \bmod m$, with two auxiliary hash functions $h_1(k)$ and
$h_2(k)$. $i$ is the number of times the hash table has been probed with key $k$. Assume two
auxiliary hash functions $h_1(k)=k \bmod 16$, $h_2(k)=1+(k \bmod 15)$, and $m=16$. Show the final
hash table content after inserting all of the following keys to an initially empty hash table in
the given order: \{16, 3, 35, 67, 51, 1, 15, 31, 19, 17\}.
\item (5\%) Prof. Alpha proposes a different double hashing method, using two auxiliary hash
functions $h_1(k)=k \bmod 16$, $h_2(k)=2(k \bmod 8)$, respectively. Assume the table size $m=16$.
Use an example to briefly explain how Prof. Alpha's method is problematic.
\end{enumerate}
\end{pastexam}

\begin{pastsoln}
\textbf{(a)} $h(k, i) = \bigl(h'(k) + i\bigr) \bmod m = \bigl((k \bmod m) + i\bigr) \bmod m$，
$i = 0, 1, \ldots, m - 1$（資結書第~12.3~節）。 最外面的 $\bmod m$ 不能省：探測走到最後一格
$m - 1$ 之後要繞回第 $0$ 格。

\textbf{(b)} 每個鍵依序試 $h(k, 0), h(k, 1), \ldots$，第一個空格就放下：
\begin{center}
\small
\begin{tabular}{rcccl}
\toprule
$k$ & $h_1(k)$ & $h_2(k)$ & 探測的格子 & 放在 \\
\midrule
16 & 0 & 2 & 0 & 0 \\
3 & 3 & 4 & 3 & 3 \\
35 & 3 & 6 & 3（有 3）, 9 & 9 \\
67 & 3 & 8 & 3, 11 & 11 \\
51 & 3 & 7 & 3, 10 & 10 \\
1 & 1 & 2 & 1 & 1 \\
15 & 15 & 1 & 15 & 15 \\
31 & 15 & 2 & 15（有 15）, 1（有 1）, 3（有 3）, 5 & 5 \\
19 & 3 & 5 & 3, 8 & 8 \\
17 & 1 & 3 & 1（有 1）, 4 & 4 \\
\bottomrule
\end{tabular}
\end{center}
例如 $31$：$h_1 = 31 \bmod 16 = 15$，$h_2 = 1 + (31 \bmod 15) = 2$，依序試
$15, (15 + 2) \bmod 16 = 1, (15 + 4) \bmod 16 = 3, (15 + 6) \bmod 16 = 5$，前三格都有鍵，放在第
$5$ 格。 最後的表（空白表示空格）：
\begin{center}
\small
\setlength{\tabcolsep}{4pt}
\begin{tabular}{|*{16}{c|}}
\hline
0 & 1 & 2 & 3 & 4 & 5 & 6 & 7 & 8 & 9 & 10 & 11 & 12 & 13 & 14 & 15 \\ \hline
16 & 1 &  & 3 & 17 & 31 &  &  & 19 & 35 & 51 & 67 &  &  &  & 15 \\ \hline
\end{tabular}
\end{center}
$35, 67, 51, 19$ 的 $h_1$ 都是 $3$，但 $h_2$ 不同，各自跳到不同的格子，不會像線性探測那樣全部擠
在第 $3$ 格後面。

\textbf{(c)} 問題在 $h_2(k) = 2(k \bmod 8)$：
\begin{enumerate}
\item \textbf{$h_2$ 可能是 $0$。} $k$ 是 $8$ 的倍數時 $h_2(k) = 0$，探測序列是
$h_1(k), h_1(k), h_1(k), \ldots$，永遠只試同一格。 例子：先插入 $24$，$h_1(24) = 8$，放在第 $8$ 格；
再插入 $8$，$h_1(8) = 8$、$h_2(8) = 0$，第 $8$ 格有鍵，之後每一次探測都還是第 $8$ 格，表裡另外
$15$ 格是空的也放不進去（插入失敗或無窮迴圈）。
\item \textbf{$h_2$ 一定是偶數}，和 $m = 16$ 不互質。 探測序列每次跳偶數格，只會走到和 $h_1(k)$
同奇偶的那 $8$ 格。 例如 $k = 3$：$h_1 = 3$、$h_2 = 6$，探測 $3, 9, 15, 5, 11, 1, 7, 13$，接著又回到
$3$，偶數格一個都碰不到。 奇數格滿了，奇數的鍵就插不進去，雖然偶數格還有空位。
\end{enumerate}
資結書定理~12.4：$\gcd(h_2(k), m) = 1$ 時探測序列才會走遍全部 $m$ 格。 $m$ 是 $2$ 的冪時，常見的
做法是讓 $h_2(k)$ 永遠是奇數，例如 $h_2(k) = 2(k \bmod 8) + 1$。 (b) 的 $h_2(k) = 1 + (k \bmod 15)$
雖然也會出現偶數（例如 $h_2(16) = 2$），但至少不會是 $0$，而且大約一半的鍵拿到奇數。
\end{pastsoln}
